\section{Estructura de la memoria}

La presente memoria se organiza en cinco capítulos y dos anexos:

El \textbf{Capítulo 2} presenta el estado del arte en producción audiovisual: la evolución
desde la producción tradicional hasta la producción remota basada en software, las
tecnologías de contribución y distribución de vídeo, y una comparativa de los
mezcladores de vídeo disponibles tanto en el ámbito hardware como en software.

El \textbf{Capítulo 3} describe el diseño e implementación del sistema de producción remota de
vídeo Voctomix 2.0: la arquitectura del sistema (voctocore y voctogui), las fuentes de
señal de cámaras y auxiliares, los modos de composición, las transiciones, el sistema de
\textit{overlays} (superposiciones) y rotulación dinámica, la telemetría mediante RabbitMQ, y el despliegue
contenerizado con Docker.

El \textbf{Capítulo 4} recoge las pruebas de validación del sistema en dos escenarios
de despliegue: el escenario de dos PCs en red local y el escenario Docker. Se presentan
mediciones de latencia, capturas de los registros de telemetría en RabbitMQ y evidencias
del correcto funcionamiento del sistema.

El \textbf{Capítulo 5} presenta las conclusiones del trabajo y las líneas de trabajo
futuro identificadas, que incluyen el despliegue en Kubernetes y la validación con
cámaras reales en el entorno del grupo CyberNemo de la UPM.

Los \textbf{Anexos} incluyen información complementaria sobre los aspectos éticos,
económicos, sociales y ambientales del proyecto (Anexo A), así como el presupuesto
económico detallado (Anexo B).
